\begin{abstract}
Este artículo presenta una revisión sistemática del estado del arte en el uso de técnicas de Inteligencia Artificial (IA) y Visión Computacional aplicadas a los Sistemas de Transporte Inteligentes (ITS) y al Monitoreo del Comportamiento de Conductores. Analizando un corpus de 213 publicaciones científicas recientes (2020--2026), clasificamos y contrastamos los enfoques metodológicos predominantes, que abarcan desde el aprendizaje profundo (Deep Learning) hasta el aprendizaje automático tradicional (Machine Learning) y técnicas de visión híbridas. Evaluamos los entornos experimentales más comunes (sistemas de videovigilancia CCTV, drones/UAV y pruebas en entornos reales), las plataformas de hardware y software empleadas para su implementación, así como las métricas de rendimiento y limitaciones reportadas en la literatura. Los resultados demuestran una clara transición hacia modelos de aprendizaje profundo en años recientes y un auge en el despliegue en hardware de borde (Edge Computing), enfrentando retos persistentes relacionados con oclusiones visuales y variaciones extremas en las condiciones de iluminación.
\end{abstract}

\begin{IEEEkeywords}
Sistemas de Transporte Inteligentes, Visión Computacional, Deep Learning, Monitoreo de Conductores, Detección de Infracciones, Aprendizaje Automático.
\end{IEEEkeywords}
